Put $s=a+b$ and $p=ab$. The equation becomes $sp+2p+s=97$, or
\[
(s+2)(p+1)=99.
\]
Thus $d=s+2$ is one of $\pm1,\pm3,\pm9,\pm11,\pm33,\pm99$, and $p=\dfrac{99}{d}-1$. The numbers $a,b$ are integer roots of $X^2-sX+p=0$, so its discriminant must be a nonnegative square with the same parity as $s$. Checking these twelve divisors gives
\[
\begin{array}{c|rrrr}
d&-9&9&11&99\\ \hline
s&-11&7&9&97\\
p&-12&10&8&0\\
\text{roots}&1,-12&2,5&1,8&0,97
\end{array}
\]
For the other divisors $d=-99,-33,-11,-3,-1,1,3,33$, the discriminants are respectively $10209,1241,209,161,409,-391,-127,953$, none a nonnegative square. Every pair in the table verifies the original equation. Hence
\[
\boxed{a+b\in\{-11,7,9,97\}}.
\]
